\section*{Abstract}
\label{sec:abstract-1}
Understanding how molecular programs are embedded within tissue morphology is a central challenge in spatial biology. While vision transformer (ViT) foundation models capture rich histological structure and spatial transcriptomics provide molecular context, existing multimodal approaches largely rely on contrastive alignment and do not directly learn joint morpho-molecular representations. We introduce an early-fusion multimodal transformer that integrates subcellular Xenium transcript readouts directly into the ViT token stream, enabling fine-grained cross-modal interaction without cell segmentation. We evaluate our approach on two complementary tasks motivated by the two-tier panel structure of the Xenium platform: (i) \emph{add-on gene prediction}, where held-out study-specific genes are predicted from the universal core panel and histology, a task directly relevant to panel optimization and augmentation; and (ii) \emph{cell type composition prediction}, a harder task that requires capturing local co-expression patterns defining cellular identity, and one of direct biological value given the laborious nature of manual annotation. Across a comprehensive benchmark of unimodal baselines and early/late-fusion variants, fusion models achieve substantial improvements in gene expression prediction and even stronger gains in cell type composition, with our early fusion approach outperforming all other architectures on the latter task. Beyond architecture comparisons, we also explore the effect of transcript representations, demonstrating that performance gains are driven primarily by spatially aligned, token-level transcript representations rather than fusion timing alone. We further show that model capacity selectively benefits multimodal fusion but not morphology-only models, implicating cross-modal interaction as the primary driver of additional representational capacity.

Panel-specific experiments further reveal that core panel genes alone are nearly sufficient for cell type prediction, while per-cell-type analysis uncovers characteristic unimodal failure modes, including a collapse of expression-only models on NK cells that is not mirrored by morphology-only models. Together, these results highlight expressive, spatially grounded fusion as a key ingredient for multimodal representation learning in spatial biology, and demonstrate that the two modalities carry complementary rather than redundant information for predicting cellular phenotypes.

